% GENERATED FILE. Edit canonical lessons, then run npm run generate:books.
% canonical-source-hash: fnv1a64:ab60f0a71ebc4b62
% canonical-lessons: HI-C311-sahayta, HI-C311-svayamsevak, HI-C311-jutna, HI-C311-jutana, HI-C311-hausla, HI-R311-review-304-311

\chapter{Lending a Hand}
\label{ch:hi-lending-a-hand}
\hlchaptermodality{311}

\begin{chapteropening}
\textbf{By the end of this chapter:} \emph{I can say who helped and what they collected.}

Complete the last lesson of chapter 311: I can say who helped and what they collected.
\end{chapteropening}

\section[\texorpdfstring{\hi{सहायता} (sahāytā)}{सहायता (sahāytā)}]{\hi{सहायता} (sahāytā) — aid; assistance}

\label{lesson:HI-C311-sahayta}



\begin{quote}
Before the new one: say the Hindi for safe and sound, then the Hindi for comfort.
\end{quote}

\subsection*{You'll want to know: \hi{सहायता}}
\textbf{\hi{सहायता}} — \emph{sahāytā} — \textquotedblleft{}aid; assistance\textquotedblright{}.

\begin{quote}
\hi{कई} \hi{देशों} \hi{ने} \hi{सहायता} \hi{भेजी।}
\end{quote}

\begin{quote}
\emph{kaī deśoṁ ne sahāytā bhejī.}
\end{quote}

\begin{quote}
Many countries sent aid.
\end{quote}

Feminine.

\begin{cousinweb}
\textbf{\hi{सहायता}} (\emph{sahāytā}) is Sanskrit \emph{sahāyatā} (Monier-Williams), from \emph{sahāya}, a companion who goes along with you. It is the formal word beside \textbf{\hi{मदद}}.
\end{cousinweb}

\begin{grammarlens}[title={What you've built}]
One more word for telling what happened.
\end{grammarlens}

\subsection*{Your turn}
\begin{itemize}
\raggedright
  \item \emph{Say it:} \emph{sahāytā}
  \item \emph{Say it:} \emph{sahāytā}, once more
  \item \emph{Say it:} \emph{kaī deśoṁ ne sahāytā bhejī.} — Many countries sent aid.
  \item \emph{Recall:} say the Hindi for safe and sound, then the Hindi for comfort, then say \emph{sahāytā} again
\end{itemize}

\begin{tcolorbox}[breakable,colback=teal!4,colframe=teal!35!black,title={Before you move on}]
What does \hi{सहायता} mean? (\textquotedblleft{}Aid; assistance\textquotedblright{}.) Say it once more.
\end{tcolorbox}
\section[\texorpdfstring{\hi{स्वयंसेवक} (svayaṁsevak)}{स्वयंसेवक (svayaṁsevak)}]{\hi{स्वयंसेवक} (svayaṁsevak) — a volunteer}

\label{lesson:HI-C311-svayamsevak}



\begin{quote}
Before the new one: say the Hindi for comfort, then the Hindi for aid.
\end{quote}

\subsection*{You'll want to know: \hi{स्वयंसेवक}}
\textbf{\hi{स्वयंसेवक}} — \emph{svayaṁsevak} — \textquotedblleft{}a volunteer\textquotedblright{}.

\begin{quote}
\hi{मुझे} \hi{उम्मीद} \hi{है} \hi{कि} \hi{कल} \hi{और} \hi{स्वयंसेवक} \hi{आएँगे।}
\end{quote}

\begin{quote}
\emph{mujhe ummīd hai ki kal aur svayaṁsevak āeṅge.}
\end{quote}

\begin{quote}
I hope more volunteers will come tomorrow.
\end{quote}

Masculine.

\begin{cousinweb}
\textbf{\hi{स्वयंसेवक}} (\emph{svayaṁsevak}) is a modern compound of Sanskrit \emph{svayam}, \textquotedblleft{}oneself, of one's own accord\textquotedblright{}, and \emph{sevaka}, \textquotedblleft{}a servant\textquotedblright{} (Monier-Williams): someone who serves because they chose to.
\end{cousinweb}

\begin{grammarlens}[title={What you've built}]
One more word for telling what happened.
\end{grammarlens}

\subsection*{Your turn}
\begin{itemize}
\raggedright
  \item \emph{Say it:} \emph{svayaṁsevak}
  \item \emph{Say it:} \emph{svayaṁsevak}, once more
  \item \emph{Say it:} \emph{mujhe ummīd hai ki kal aur svayaṁsevak āeṅge.} — I hope more volunteers will come tomorrow.
  \item \emph{Recall:} say the Hindi for comfort, then the Hindi for aid, then say \emph{svayaṁsevak} again
\end{itemize}

\begin{tcolorbox}[breakable,colback=teal!4,colframe=teal!35!black,title={Before you move on}]
What does \hi{स्वयंसेवक} mean? (\textquotedblleft{}A volunteer\textquotedblright{}.) Say it once more.
\end{tcolorbox}
\section[\texorpdfstring{\hi{जुटना} (juṭnā)}{जुटना (juṭnā)}]{\hi{जुटना} (juṭnā) — to gather; to come together}

\label{lesson:HI-C311-jutna}



\begin{quote}
Before the new one: say the Hindi for aid, then the Hindi for a volunteer.
\end{quote}

\subsection*{You'll want to know: \hi{जुटना}}
\textbf{\hi{जुटना}} — \emph{juṭnā} — \textquotedblleft{}to gather; to come together\textquotedblright{}.

\begin{quote}
\hi{मदद} \hi{के} \hi{लिए} \hi{स्कूल} \hi{में} \hi{सौ} \hi{लोग} \hi{जुटे।}
\end{quote}

\begin{quote}
\emph{madad ke liye skūl meṁ sau log juṭe.}
\end{quote}

\begin{quote}
A hundred people gathered at the school to help.
\end{quote}

Past (it …): \hi{वह} \hi{भीड़} \hi{जुटी}, \emph{vah bhīṛ juṭī}.

\begin{cousinweb}
\textbf{\hi{जुटना}} (\emph{juṭnā}) is inherited from \emph{*yuṭati}, \textquotedblleft{}is joined\textquotedblright{}, a form CDIAL reconstructs (CDIAL 10496), the family of \textbf{\hi{जोड़ना}}, to join. People \textbf{\hi{जुटते} \hi{हैं}} when they come together for one purpose.
\end{cousinweb}

\begin{grammarlens}[title={What you've built}]
One more word for telling what happened.
\end{grammarlens}

\subsection*{Your turn}
\begin{itemize}
\raggedright
  \item \emph{Say it:} \emph{juṭnā}
  \item \emph{Say it:} \emph{juṭnā}, once more
  \item \emph{Say it:} \emph{madad ke liye skūl meṁ sau log juṭe.} — A hundred people gathered at the school to help.
  \item \emph{Recall:} say the Hindi for aid, then the Hindi for a volunteer, then say \emph{juṭnā} again
\end{itemize}

\begin{tcolorbox}[breakable,colback=teal!4,colframe=teal!35!black,title={Before you move on}]
What does \hi{जुटना} mean? (\textquotedblleft{}To gather; to come together\textquotedblright{}.) Say it once more.
\end{tcolorbox}
\section[\texorpdfstring{\hi{जुटाना} (juṭānā)}{जुटाना (juṭānā)}]{\hi{जुटाना} (juṭānā) — to collect; to raise (money)}

\label{lesson:HI-C311-jutana}



\begin{quote}
Before the new one: say the Hindi for a volunteer, then the Hindi for to gather.
\end{quote}

\subsection*{You'll want to know: \hi{जुटाना}}
\textbf{\hi{जुटाना}} — \emph{juṭānā} — \textquotedblleft{}to collect; to raise (money)\textquotedblright{}.

\begin{quote}
\hi{स्वयंसेवकों} \hi{ने} \hi{एक} \hi{लाख} \hi{रुपये} \hi{जुटाए।}
\end{quote}

\begin{quote}
\emph{svayaṁsevakoṁ ne ek lākh rupaye juṭāe.}
\end{quote}

\begin{quote}
The volunteers raised a hundred thousand rupees.
\end{quote}

Past, with \hi{ने} (he or she …): \hi{उसने} \hi{पैसा} \hi{जुटाया}, \emph{usne paisā juṭāyā}.

\begin{cousinweb}
\textbf{\hi{जुटाना}} (\emph{juṭānā}) is the causative of \textbf{\hi{जुटना}}, to gather, from the same \emph{*yuṭati} (CDIAL 10496): to make things come together. You \textbf{\hi{जुटाते} \hi{हो}} money, help or courage.
\end{cousinweb}

\begin{grammarlens}[title={What you've built}]
One more word for telling what happened.
\end{grammarlens}

\subsection*{Your turn}
\begin{itemize}
\raggedright
  \item \emph{Say it:} \emph{juṭānā}
  \item \emph{Say it:} \emph{juṭānā}, once more
  \item \emph{Say it:} \emph{svayaṁsevakoṁ ne ek lākh rupaye juṭāe.} — The volunteers raised a hundred thousand rupees.
  \item \emph{Recall:} say the Hindi for a volunteer, then the Hindi for to gather, then say \emph{juṭānā} again
\end{itemize}

\begin{tcolorbox}[breakable,colback=teal!4,colframe=teal!35!black,title={Before you move on}]
What does \hi{जुटाना} mean? (\textquotedblleft{}To collect; to raise (money)\textquotedblright{}.) Say it once more.
\end{tcolorbox}
\section[\texorpdfstring{\hi{हौसला} (hauslā)}{हौसला (hauslā)}]{\hi{हौसला} (hauslā) — spirit; morale}

\label{lesson:HI-C311-hausla}



\begin{quote}
Before the new one: say the Hindi for to gather, then the Hindi for to collect.
\end{quote}

\subsection*{You'll want to know: \hi{हौसला}}
\textbf{\hi{हौसला}} — \emph{hauslā} — \textquotedblleft{}spirit; morale\textquotedblright{}.

\begin{quote}
\hi{स्वयंसेवकों} \hi{ने} \hi{पीड़ितों} \hi{का} \hi{हौसला} \hi{बढ़ाया।}
\end{quote}

\begin{quote}
\emph{svayaṁsevakoṁ ne pīṛitoṁ kā hauslā baṛhāyā.}
\end{quote}

\begin{quote}
The volunteers raised the victims' spirits.
\end{quote}

Masculine.

\begin{cousinweb}
\textbf{\hi{हौसला}} (\emph{hauslā}) is usually traced, through Persian, to Arabic \emph{ḥawṣala}, first the crop of a bird and then capacity; in Hindi it is the courage to keep going. \textbf{\hi{हौसला} \hi{बढ़ाना}} is to encourage.
\end{cousinweb}

\begin{grammarlens}[title={What you've built}]
One more word for telling what happened.
\end{grammarlens}

\subsection*{Your turn}
\begin{itemize}
\raggedright
  \item \emph{Say it:} \emph{hauslā}
  \item \emph{Say it:} \emph{hauslā}, once more
  \item \emph{Say it:} \emph{svayaṁsevakoṁ ne pīṛitoṁ kā hauslā baṛhāyā.} — The volunteers raised the victims' spirits.
  \item \emph{Recall:} say the Hindi for to gather, then the Hindi for to collect, then say \emph{hauslā} again
\end{itemize}

\begin{tcolorbox}[breakable,colback=teal!4,colframe=teal!35!black,title={Before you move on}]
What does \hi{हौसला} mean? (\textquotedblleft{}Spirit; morale\textquotedblright{}.) Say it once more.
\end{tcolorbox}
\section[(dialogue)]{Looking back: from The Storm to Lending a Hand}

\label{lesson:HI-R311-review-304-311}



\begin{quote}
Say the Hindi for darkness and to thunder.
\end{quote}

\subsection*{Your turn}
Say these aloud:

\begin{itemize}
\raggedright
  \item darkness, to thunder, to flash, torrential (of rain), destruction
  \item to swell up, a dam, little by little, a deluge, a whirlpool
  \item a famine, starvation, to rot, an epidemic, cholera
  \item to be trapped, to call out (to), a dead silence, to take out, in a sorry state
  \item a raft, a diver, alive, to set free, to dart
  \item relief, a camp, an encampment, a victim, refuge
  \item to be separated, to wander, a search, safe and sound, comfort
  \item aid, a volunteer, to gather, to collect, spirit
\end{itemize}

\begin{tcolorbox}[breakable,colback=teal!4,colframe=teal!35!black,title={Before you move on}]
Darkness? (\textbf{\hi{अँधेरा}}, \emph{ãdherā}.) To thunder? (\textbf{\hi{गरजना}}, \emph{garajnā}.) To flash? (\textbf{\hi{चमकना}}, \emph{camaknā}.) Torrential (of rain)? (\textbf{\hi{मूसलाधार}}, \emph{mūslādhār}.) Destruction? (\textbf{\hi{तबाही}}, \emph{tabāhī}.) To swell up? (\textbf{\hi{उफनना}}, \emph{uphannā}.) A dam? (\textbf{\hi{बाँध}}, \emph{bā̃dh}.) Little by little? (\textbf{\hi{क्रमशः}}, \emph{kramśaḥ}.) A deluge? (\textbf{\hi{सैलाब}}, \emph{sailāb}.) A whirlpool? (\textbf{\hi{भँवर}}, \emph{bhãvar}.) A famine? (\textbf{\hi{अकाल}}, \emph{akāl}.) Starvation? (\textbf{\hi{भुखमरी}}, \emph{bhukhmarī}.) To rot? (\textbf{\hi{सड़ना}}, \emph{saṛnā}.) An epidemic? (\textbf{\hi{महामारी}}, \emph{mahāmārī}.) Cholera? (\textbf{\hi{हैज़ा}}, \emph{haizā}.) To be trapped? (\textbf{\hi{फँसना}}, \emph{phãsnā}.) To call out (to)? (\textbf{\hi{पुकारना}}, \emph{pukārnā}.) A dead silence? (\textbf{\hi{सन्नाटा}}, \emph{sannāṭā}.) To take out? (\textbf{\hi{निकालना}}, \emph{nikālnā}.) In a sorry state? (\textbf{\hi{बेहाल}}, \emph{behāl}.) A raft? (\textbf{\hi{बेड़ा}}, \emph{beṛā}.) A diver? (\textbf{\hi{गोताख़ोर}}, \emph{gotākhor}.) Alive? (\textbf{\hi{ज़िंदा}}, \emph{zindā}.) To set free? (\textbf{\hi{छुड़ाना}}, \emph{chuṛānā}.) To dart? (\textbf{\hi{लपकना}}, \emph{lapaknā}.) Relief? (\textbf{\hi{राहत}}, \emph{rāhat}.) A camp? (\textbf{\hi{शिविर}}, \emph{śivir}.) An encampment? (\textbf{\hi{डेरा}}, \emph{ḍerā}.) A victim? (\textbf{\hi{पीड़ित}}, \emph{pīṛit}.) Refuge? (\textbf{\hi{शरण}}, \emph{śaraṇ}.) To be separated? (\textbf{\hi{बिछुड़ना}}, \emph{bichuṛnā}.) To wander? (\textbf{\hi{भटकना}}, \emph{bhaṭaknā}.) A search? (\textbf{\hi{तलाश}}, \emph{talāś}.) Safe and sound? (\textbf{\hi{सलामत}}, \emph{salāmat}.) Comfort? (\textbf{\hi{दिलासा}}, \emph{dilāsā}.) Aid? (\textbf{\hi{सहायता}}, \emph{sahāytā}.) A volunteer? (\textbf{\hi{स्वयंसेवक}}, \emph{svayaṁsevak}.) To gather? (\textbf{\hi{जुटना}}, \emph{juṭnā}.) To collect? (\textbf{\hi{जुटाना}}, \emph{juṭānā}.) Spirit? (\textbf{\hi{हौसला}}, \emph{hauslā}.)
\end{tcolorbox}
